\section{Integral relation lattices}
\label{sec:integral}

Let
\[
 L_i=\langle\,T:T\in\cC_i\,\rangle_{\ZZ}
 \subseteq K_i^{\ZZ},
 \qquad
 \Theta_i=K_i^{\ZZ}/L_i.
\]
The rational theorem determines the rank but not the index of \(L_i\).
This distinction is the source of the integral phenomenon below.

\begin{lemma}[Finiteness and local saturation]
\label{lem:local-saturation}
The group \(\Theta_i\) is finite.  Let \(p\) be a prime.  If the image of
\(L_i\otimes\ZZ_p\) spans
\((K_i^{\ZZ}\otimes\ZZ_p)/p(K_i^{\ZZ}\otimes\ZZ_p)\), then
\(\Theta_i[p^\infty]=0\).
\end{lemma}

\begin{proof}
The lattice \(K_i^{\ZZ}\) is primitive in \(\ZZ[\cB_i]\): its quotient
is the image of \(A_i\), a subgroup of the free group \(\ZZ[\cP]\).
By \cref{thm:full-kernel}, \(L_i\) and \(K_i^{\ZZ}\) have the same
rational span and hence the same rank.  Their quotient is finite.

Put \(M=K_i^{\ZZ}\otimes\ZZ_p\) and
\(N=L_i\otimes\ZZ_p\).  The spanning hypothesis is \(N+pM=M\).
Nakayama's lemma gives \(N=M\), so the localized quotient vanishes.
\end{proof}

\subsection{The prime \(2\)}

\begin{proposition}[Binary saturation]
\label{prop:p2}
For \(i=1,2,3\),
\[
  \Theta_i[2^\infty]=0.
\]
\end{proposition}

\begin{proof}
In each design, an explicit set of \(66\,674\) minimum trades is linearly
independent in \(\FF_2[\cB_i]\).  Primitivity of \(K_i^{\ZZ}\) makes
\(K_i^{\ZZ}/2K_i^{\ZZ}\) a \(66\,674\)-dimensional subspace of that
ambient permutation module.  The selected trades lie in this subspace and
already have its full dimension.  Apply \cref{lem:local-saturation}.
\end{proof}

This statement concerns liftable integral relations, not the complete
modular null-design space.  The distinction is especially visible at the
next prime and is consistent with the modular minimum-support setting of
\cite{Krotov2020}.

\subsection{The prime \(3\): the liftable kernel}

Let \(\alpha:\ZZ[\cB_i]\to\ZZ\) be augmentation,
\(\alpha(\sum n_BB)=\sum n_B\).  Summing all point coordinates of an
integral incidence relation gives
\[
 15\,\alpha(x)=0.
\]
Thus \(\alpha\) vanishes on \(K_i^{\ZZ}\), even though its reduction
modulo \(3\) is not a consequence of the incidence equations.

\begin{lemma}[Augmentation--Bockstein identification]
\label{lem:p3-lift}
For every \(i\),
\[
 K_i^{\ZZ}/3K_i^{\ZZ}
 =
 \ker_{\FF_3}
 \begin{bmatrix}A_i\\ \alpha\end{bmatrix}.
\]
\end{lemma}

\begin{proof}
The locked Smith data have exactly one factor divisible by \(3\), namely
\(120\), in the \(255\)-row incidence presentation.  In particular,
\[
 \rank_{\FF_3}A_i=254,\qquad
 \rank_{\FF_3}\begin{bmatrix}A_i\\ \alpha\end{bmatrix}=255.
\]
The modular kernel of \(A_i\) therefore has dimension \(66\,675\), one
larger than the reduction of the primitive rank-\(66\,674\) integral
kernel.  The augmented kernel has dimension \(66\,674\) and contains that
reduction by the preceding augmentation identity.  Equality follows.
Equivalently, the missing modular direction is the one-dimensional
\(\Tor_1\) contribution produced by the factor \(120\).
\end{proof}

Write \(\GammaG=F\times C_3\), where \(|F|=68\), and put
\[
 R_3=\FF_3[C_3]\cong\FF_3[u]/(u^3),
 \qquad u=c-1.
\]
The ideal \((u)\) is the radical of \(\FF_3\GammaG\), while
\(\FF_3F\) is semisimple.  Let \(M_i^{(3)}\) denote the liftable kernel in
\cref{lem:p3-lift}, and let \(\overline M_i^{(3)}=M_i^{(3)}/uM_i^{(3)}\)
be its semisimple top.

The exact component certificates decompose this top into four Frobenius
families, whose Wedderburn dimensions account for all \(68\) dimensions
of \(\FF_3F\).  The images of minimum trades attain total top dimensions
\[
\begin{array}{c|rrr}
 &\cB_1&\cB_2&\cB_3\\
\midrule
\dim_{\FF_3}\overline M_i^{(3)}
 &22\,304&22\,338&22\,270.
\end{array}
\]
Each component rank is accompanied by a maximal minor and a rank-loss
control.

\begin{proposition}[Ternary saturation]
\label{prop:p3}
For \(i=1,2,3\),
\[
  \Theta_i[3^\infty]=0.
\]
\end{proposition}

\begin{proof}
By \cref{lem:p3-lift}, the module tested by the certificates is exactly
the reduction of the integral kernel, not the one-dimension-larger
modular kernel.  The component ranks say that the minimum trades generate
every simple summand of \(\overline M_i^{(3)}\).  Hence
\[
  M_i^{(3)}=L_i^{(3)}+uM_i^{(3)}.
\]
The ideal \((u)\) is nilpotent, so Nakayama's lemma gives
\(M_i^{(3)}=L_i^{(3)}\).  Applying the local criterion to the integral
lattices proves the assertion.
\end{proof}

\subsection{The prime \(17\): a free cyclic kernel}

The normal subgroup \(C_{17}\) acts freely on points and blocks.  There
are \(15\) point orbits and \(3\,937\) block orbits.  Put
\[
 R_{17}=\FF_{17}[C_{17}]
       \cong\FF_{17}[u]/(u^{17}),\qquad u=a-1.
\]
The two permutation modules are free \(R_{17}\)-modules of ranks
\(3\,937\) and \(15\).  The coinvariant incidence map is surjective.
Nakayama therefore makes the full \(R_{17}\)-map surjective, and its
kernel \(M_i^{(17)}\) is free of rank
\[
 3\,937-15=3\,922.
\]

The quotient \(H=\GammaG/C_{17}\cong C_4\times C_3\) is semisimple in
characteristic \(17\).  Its eight simple components have total Wedderburn
mass \(12=|H|\).  In every component, exact fixed-space matrices show
that the minimum trades span
\[
 M_i^{(17)}/uM_i^{(17)}.
\]

\begin{proposition}[Seventeen-adic saturation]
\label{prop:p17}
For \(i=1,2,3\),
\[
  \Theta_i[17^\infty]=0.
\]
\end{proposition}

\begin{proof}
The eight component ranks generate the complete semisimple top of the free
\(R_{17}\)-kernel.  Thus
\(M_i^{(17)}=L_i^{(17)}+uM_i^{(17)}\).  Nakayama over the local ring
\(R_{17}\) gives equality, and \cref{lem:local-saturation} lifts the
conclusion to the \(17\)-adic lattices.
\end{proof}

\subsection{What the small components exclude}

Away from \(2\cdot3\cdot17\), the six rational types of dimensions at
most four admit integral component models.  For each of the resulting
\(18\) maps, two exact maximal minors have the gcd shown below.
\[
\begin{array}{c|rrr}
V&\cB_1&\cB_2&\cB_3\\
\midrule
\mathbf1&24&6&24\\
\eps&8&2&128\\
\iot&1024&4096&128\\
\om&4&216&8\\
\eps\om&108&12&2\\
\iot\om&33\,554\,432&16\,777\,216&49\,152
\end{array}
\]
Every displayed number is supported on \(\{2,3,17\}\).  Bezout
combination of the corresponding minors makes each maximal-minor ideal a
unit after these primes are inverted.

\begin{proposition}[Small-component exclusion]
\label{prop:small-components}
No prime outside \(\{2,3,17\}\) occurs in any of the \(18\) component
cokernels belonging to
\[
\mathbf1,\eps,\iot,\om,\eps\om,\iot\om.
\]
\end{proposition}

This is not a support theorem for \(\Theta_i\).  The six large maps
\(\rh/\cB_i\) and \(\rh\om/\cB_i\) require separate complete-equation-map
analysis.  A control-prime-selected presentation can display a good-prime
rank defect that disappears when the full minimum-circuit space is imposed.
\Cref{sec:137} resolves the selected primes \(137\), \(219\,097\), and
\(3\,288\,036\,131\); \cref{sec:global-rho} then closes the single
\(\rh/\cB_1\) component globally away from \(2\), \(3\), and \(17\).
